

We presented \MethodName{}, a general-purpose robot foundation model that exhibits out-of-the-box compositional generalization, effective language following, and task performance that is competitive with more specialized models that are fine-tuned to individual dexterous tasks. At the core of \MethodName{} is a diverse prompting strategy, inspired by prompt expansion, where additional information about the episode is provided to the model during training and, optionally, at test time. This additional information includes more detailed language, episode metadata, and subgoal images. Our experiments show that by using diverse prompting and a larger and more diverse dataset, \MethodName{} can represent policies of many different qualities, and distill specialist performance back into one pre-trained model. \MethodName{}  acquires a number of emergent capabilities, such as the ability to transfer skills across robots, follow complex language commands, and generalize compositionally, recombining skills in new ways to solve new tasks.

While \MethodName{} generalizes broadly, the success rate for zero-shot generalization is (unsurprisingly) lower than in-distribution tasks: while seen tasks often have success rates in excess of 90\%, unseen tasks or unseen task-robot combinations have success rates in the 60-80\% range. An exciting direction for future work is to leverage the high steerability of \MethodName{} to efficiently learn from data in the test task, for example with more detailed language coaching or even with autonomous reinforcement learning.

A perhaps surprising limitation of our experiments is that it's practically difficult when training on such large and diverse datasets to definitively determine which tasks are truly ``seen'' or ``unseen'': while some of our generalization experiments (e.g., those in Sec.~\ref{sec:compositional_generalization}) use tasks for which we did not deliberately collect data, our dataset contains so many different scenes and behaviors that potentially related skills may well be present elsewhere in the data, either with a different label, or incidentally as part of performing other tasks. In many ways this mirrors the challenge of understanding generalization with large language models: determining what is truly novel becomes difficult, and the model may well be achieving generalization primarily by ``remixing'' skills and behaviors from other situations. However, we would posit that this in fact is the essence of \emph{compositional generalization}. Practically, whether the behaviors are truly new or merely novel combinations of seen parts, the ramifications are similar: instead of deliberating collecting targeted data for each new task that we want the robot to solve, a model that provides compositional generalization would allow the user to simply \emph{prompt} it to do the desired task. Models that can enable such compositional generalization at scale would transform how we approach robotic learning, making it possible to prompt, coach, and explain things to a robot rather than needing to collect additional action data.

